\section{Conclusões}

Esta atividade laboratorial teve como objetivo aplicar corretamente as funções de voltímetro e amperímetro do multímetro e verificar experimentalmente as Leis de Kirchhoff, nomeadamente a Lei das Malhas, bem como o Princípio da Sobreposição.

\subsection*{Verificação da Lei das Malhas}

A Lei das Malhas estabelece que a soma algébrica das diferenças de potencial numa malha fechada é nula. A partir das medições realizadas, obteve-se para a primeira malha uma soma de tensões igual a $(0,051 \pm 0,03)\,\text{V}$ e, para a segunda malha, um valor praticamente nulo dentro da incerteza experimental.

Embora o valor ideal previsto teoricamente seja exatamente zero, o desvio observado é muito reduzido quando comparado com as tensões individuais presentes no circuito, representando um erro relativo baixo. Assim, considerando a incerteza associada às medições, os resultados obtidos são compatíveis com a Lei das Malhas.

As pequenas discrepâncias verificadas podem ser justificadas por fatores experimentais, tais como a resolução limitada do multímetro ($\pm 0,01$ V), as tolerâncias nominais das resistências utilizadas, a presença de resistências internas nas fontes e nos condutores, bem como possíveis quedas de tensão adicionais introduzidas pelos próprios instrumentos de medição.

\subsection*{Verificação do Princípio da Sobreposição}

Na análise do princípio da sobreposição, avaliou-se a contribuição individual de cada fonte para a tensão $V_2$. A soma das tensões medidas com cada fonte ativa separadamente resultou em
\[
    V_2^{15V} + V_2^{7.5V} = (14,87 \pm 0,02)\,\text{V},
\]
enquanto a medição com ambas as fontes ligadas simultaneamente forneceu
\[
    V_2 = (14,86 \pm 0,01)\,\text{V}.
\]

A diferença entre estes valores é inferior à incerteza experimental, pelo que se pode considerar que, para esta grandeza, o princípio da sobreposição foi validado experimentalmente.

Contudo, para a corrente $I_2$, verificou-se que o valor experimental não se encontra dentro do intervalo esperado a partir da soma das contribuições individuais. Esta divergência poderá estar associada à influência da resistência interna do amperímetro, à perturbação introduzida pela sua inserção em série no circuito e à acumulação de erros sistemáticos não considerados no modelo teórico ideal.

\subsection*{Avaliação da Precisão e Exatidão}

Em termos de precisão, as medições apresentaram baixa dispersão e uma incerteza reduzida ($\pm 0,01$ V), indicando boa consistência dos valores obtidos.

Relativamente à exatidão, os resultados mostram uma concordância global satisfatória com os valores teóricos previstos, particularmente na verificação da Lei das Malhas e da tensão $V_2$ pelo princípio da sobreposição. As diferenças observadas são compatíveis com as limitações inerentes ao processo experimental e não comprometem a validade das leis físicas analisadas.

\bigskip

Conclui-se, assim, que os objetivos da experiência foram alcançados. Foi possível utilizar corretamente o multímetro nas suas diferentes configurações e confirmar, dentro das limitações experimentais, a validade da Lei das Malhas e do princípio da sobreposição no circuito estudado.